
Distribution-aware dataset search addresses the challenge of discovering datasets by statistical property. In a setting where the raw data is not accessible to the searcher (e.g., due to privacy or storage constraints), only per-column histograms are available. Fainder~\cite{behme_2024} introduces an index for this problem: it clusters histograms with K-means, aligns them onto a shared bin grid inside each, and precomputes the cumulative-density arrays so that percentile predicates are answered by binary search, rather than histogram scans. Using the GitTables corpus of 5M histograms, it achieves a speedup of more than two orders of magnitude over a na\"{i}ve profile-scan baseline, shrinking the index from over 1\,TB to 3.2\,GB. Fainder was originally written in Python, and the Python overhead is the dominant factor in its query time: the GIL (Global Interpreter Lock) serialises the per-cluster loop across threads, its AoS (Array-of-Structs) layout pollutes cache lines with unread hist-id fields, and per-iteration interpreter dispatch compounds onto both. This results in a $2.5$-hour runtime for 10{,}000 queries on the full five-million-histogram workload, regardless of thread count.

This thesis is a re-implementation of Fainder's query phase in Rust, a systems programming language with a strong focus on performance and memory safety. The engine resulting from this re-implementation is used to characterise the performance ceilings of distribution-aware search on modern multi-core hardware, and to investigate the composability of optimisations that address these ceilings. We find that distribution-aware search is bounded by four distinct performance ceilings: single-thread compute / L1 latency, shared L3 capacity, on-socket memory pressure, and cross-socket UPI interconnect. Each binds at a different thread-count and has its own class of intervention. Another finding is the ``negative composability'' of optimisations: independent ablations on five different axes all displaying the same behaviour, a finer optimisation layering on top of a coarser that has already reduced the binding ceiling, and the finer optimisation's overhead emerging as a net regression rather than a diminishing return.

At the scaling cell $t=16$ (the peak of the shared-memory scaling curve), the Rust engine delivers a $530$--$598\times$ end-to-end speedup over the Python baseline on precision- and recall-preserving executions. This ratio is due to Rust's avoidance of per-\emph{(query, cluster)} result-dict construction in Python, leaving the search phase with a $22.7\times$ speedup on $324$K histograms, dropping to $1.34\times$ on $5.02$M as the workload approaches the shared hardware ceilings this thesis characterises. Further optimisations include single-socket NUMA placement, which shaves a further $11$--$13\%$ off the wall at $t \leq 48$; the Rust port of the rebinning kernel, with a $\sim\!4.7\times$ speedup on the alignment phase of index construction for datasets under one million histograms, dropping to $\sim\!2.2\times$ at five million (due to the kernel's parallelism ceiling of ${\sim}2$ effective cores on the shared per-cluster allocation). The findings are consolidated into a dispatch policy that selects the best build per workload regime, validated within $5\%$ on $23$ of $24$ cells including an out-of-distribution density point. We consider these figures a downstream consequence of the ceiling-aware re-engineering, not the thesis's principal claim.

The methodological contribution of this thesis is the \emph{pre-flight ceiling identification}: before choosing an optimisation, we probe the target composition with hardware performance counters and rule out any optimisation whose intended axis no longer binds there. This pattern is expected to hold for any similar percentile-predicate query engine with an access pattern that is a chain of per-cluster binary searches.
